\scanpage{25}
\mspage{10}

(3.1), $c_4=0$ si et seulement si $a_2^2=0$. Supposons qu'on ait même $a_2=0$
\[
\begin{aligned}[t]
 &y^2=x^3+a_4 x+a_6 \qquad \pi=a_4^{-1}\,dy=dx/2y\\[.65em]
 &c_4=c_6=0 \qquad \Delta=-a_4^3 \, .
\end{aligned}\tag*{(5.1)}
\]

Les autres systèmes $(\pi,x,y)$ donnant une équation de cette forme s'obtiennent en faisant $s=t=0$ dans (1.2) et
\[
\begin{aligned}
u^4 a'_4 &= a_4\\
u^6 a'_6 &= a_6+r\,a_4+r^3 \, .
\end{aligned}\tag*{(5.2)}
\]

b) $p=2$\mbox{\begingroup\edef\encoutersep{\the\dimexpr.5\pgflinewidth\relax}\setbox0=\hbox{.}\edef\encwda{\the\wd0}\edef\enchta{\the\ht0}\edef\encdpa{\the\dp0}\rlap{\smash{\tikz[baseline=0pt]{\begin{scope}[shift={(\the\dimexpr \encwda/2\relax,\the\dimexpr -\encdpa/2+\enchta/2\relax)}]\handwritten{\draw[thin,gray] (0.13em,1.1em) .. controls (-0.42em,1.25em) and (-0.57em,-1.52em) .. (-0.01em,-1.45em) .. controls (0.40em,-1.46em) and (0.77em,0.33em) .. (0.75em,1.44em);}\end{scope}\pgfresetboundingbox\path[use as bounding box] (0pt,\the\dimexpr-\encdpa\relax) rectangle (\encwda,\enchta);}}}\copy0\endgroup}

En caractéristique 2 , on a $c_4=a_1^4$. Supposons qu'on ait $a_1=0$ ; $\pi$ étant donnée, on peut alors choisir $x$ et $y$ tels que (1.3) soit
\[
\begin{aligned}[t]
 &y^2+a_3\,y=x^3+a_4\,x+a_6 \qquad
\pi=\frac{dx}{a_3}=dy/\mbox{\begingroup\edef\encoutersep{\the\dimexpr.5\pgflinewidth\relax}\setbox0=\hbox{$x^2+a_4$}\edef\encwda{\the\wd0}\edef\enchta{\the\ht0}\edef\encdpa{\the\dp0}\rlap{\smash{\tikz[baseline=0pt]{\begin{scope}[shift={(\the\dimexpr \encwda/2\relax,\the\dimexpr -\encdpa/2+\enchta/2\relax)}]\coordinate (enc-den-north-east) at (\the\dimexpr \encoutersep+\encwda/2\relax,\the\dimexpr \encdpa/2+\enchta/2+\encoutersep\relax);\coordinate (enc-den-north-west) at (\the\dimexpr -\encoutersep-\encwda/2\relax,\the\dimexpr \encdpa/2+\enchta/2+\encoutersep\relax);\coordinate (enc-den-south-east) at (\the\dimexpr \encoutersep+\encwda/2\relax,\the\dimexpr -\encdpa/2-\enchta/2-\encoutersep\relax);\coordinate (enc-den-south-west) at (\the\dimexpr -\encoutersep-\encwda/2\relax,\the\dimexpr -\encdpa/2-\enchta/2-\encoutersep\relax);\handwritten{\draw[<->,thin] ([yshift=0.6ex]enc-den-north-west) .. controls ([xshift=-0.3em,yshift=-1.1ex]enc-den-south-west) and ([xshift=-0.1em,yshift=-1.1ex]enc-den-south-east) .. ([xshift=0.32em,yshift=-0.35ex]enc-den-south-east) .. controls ([xshift=0.6em]enc-den-south-east) and ([xshift=0.32em]enc-den-north-east) .. ([xshift=0.32em,yshift=0.65ex]enc-den-north-east);}\end{scope}\pgfresetboundingbox\path[use as bounding box] (0pt,\the\dimexpr-\encdpa\relax) rectangle (\encwda,\enchta);}}}\copy0\endgroup}
\\[.65em]
 &b_2=b_4=0 \qquad b_6=a_3^2 \qquad b_8=a_4^2 \qquad
\Delta=a_3^4 \qquad c_4=c_6=0 \, .
\end{aligned}\tag*{(5.3)}
\]

Les autres systèmes $(\mbox{\begingroup\edef\encoutersep{\the\dimexpr.5\pgflinewidth\relax}\setbox0=\hbox{$\pi,$}\edef\encwda{\the\wd0}\edef\enchta{\the\ht0}\edef\encdpa{\the\dp0}\rlap{\smash{\tikz[baseline=0pt]{\begin{scope}[shift={(\the\dimexpr \encwda/2\relax,\the\dimexpr -\encdpa/2+\enchta/2\relax)}]\coordinate (enc-pc-center) at (0pt,0pt);\handwritten{\draw[thin,gray,rotate around={-9:(enc-pc-center)}] (enc-pc-center) ellipse (0.62em and 1.07em);}\end{scope}\pgfresetboundingbox\path[use as bounding box] (0pt,\the\dimexpr-\encdpa\relax) rectangle (\encwda,\enchta);}}}\copy0\endgroup}\,x,y)$ donnant une équation de cette forme s'obtiennent en faisant $r=s^2$ dans (1.2) et
\[
\begin{aligned}
u^3 a'_3 &= a_3\\
u^4 a'_4 &= a_4+sa_3+s^4\\
u^6 a'_6 &= a_6+s^2\,a_4+ta_3+s^6+t^2 \, .
\end{aligned}\tag*{(5.4)}
\]
% archive sheet number: 24
